\documentclass[12pt, a4paper]{article}
\usepackage[portuguese]{babel}
\usepackage{amsmath}
\usepackage{amssymb}
\usepackage{amsfonts}
\usepackage{graphicx}
\usepackage[left=3cm,right=3cm,top=2.5cm,bottom=2.5cm]{geometry}
\usepackage{subcaption} 
\usepackage{hyperref}
\usepackage{booktabs}
\usepackage[ruled,vlined,linesnumbered]{algorithm2e}
\usepackage{float}

% Configurações de termos personalizados em português para o algorithm2e
\SetKwInput{Input}{Entrada}
\SetKwInput{Output}{Saída}
\SetKwIF{If}{ElseIf}{Else}{se}{então}{senão se}{senão}{fim se}
\SetKwFor{Para}{para}{faça}{fim para}
\SetKw{Retorna}{retorne}

\def\BibTeX{{\rm B\kern-.05em{\sc i\kern-.025em b}\kern-.08em
    T\kern-.1667em\lower.7ex\hbox{E}\kern-.125emX}}

\begin{document}

\title{\textbf{Trabalho Prático de Aprendizado de Máquina Via Álgebra Linear}}

\author{\textbf{Déborah Brito Yamamoto} \\ 
Matrícula: 2022069409 \\
Departamento de Ciência da Computação \\
Universidade Federal de Minas Gerais (UFMG)}
\date{\today}
\maketitle

\section{Introdução}

O treinamento de modelos modernos de aprendizado profundo consiste, essencialmente, em um grande problema de otimização numérica. Dado um conjunto de exemplos rotulados, deseja-se encontrar um conjunto de parâmetros capazes de minimizar uma função de perda (\textit{loss function}), isto é, uma medida quantitativa do erro cometido pelo modelo durante suas predições.

Matematicamente, seja

\[
\theta \in \mathbb{R}^{n}
\]

o vetor que reúne todos os parâmetros treináveis da rede neural, onde $n$ pode facilmente atingir dezenas ou centenas de milhões em arquiteturas modernas. Define-se uma função de perda

\[
\mathcal{L}(\theta),
\]

que associa um valor escalar à qualidade da solução atual. Quanto menor esse valor, melhor o desempenho do modelo sobre os dados de treinamento.

O problema de treinamento pode então ser escrito como

\[
\theta^{*}
=
\arg\min_{\theta}
\mathcal{L}(\theta).
\]

À primeira vista, esse objetivo parece suficiente: basta encontrar o menor valor possível da função de perda. Entretanto, a prática demonstra que essa estratégia é incompleta. Em problemas reais, o objetivo não consiste em memorizar os exemplos presentes no conjunto de treinamento, mas sim aprender padrões capazes de serem utilizados corretamente em exemplos nunca observados anteriormente.

Essa propriedade recebe o nome de \textbf{generalização}.

Em outras palavras, um modelo apresenta boa capacidade de generalização quando mantém um desempenho elevado mesmo sobre dados provenientes da mesma distribuição estatística, mas que não participaram do processo de treinamento.

\subsection{Generalização e Overfitting}

Considere dois modelos distintos treinados sobre exatamente o mesmo conjunto de dados.

Ambos podem alcançar erro praticamente nulo durante o treinamento. No entanto, quando avaliados sobre novos exemplos, um deles pode apresentar desempenho significativamente inferior.

Esse fenômeno é conhecido como \textbf{overfitting}. Nele, o modelo aprende detalhes extremamente específicos dos exemplos de treinamento — incluindo ruídos e particularidades irrelevantes — em vez de capturar regularidades gerais da distribuição dos dados.

Esse comportamento mostra que minimizar apenas a perda empírica

\[
\mathcal{L}_{train}
\]

não garante, necessariamente, uma pequena perda sobre o conjunto de teste

\[
\mathcal{L}_{test}.
\]

Consequentemente, surge uma questão fundamental:

\begin{center}
\textit{Por que dois modelos com exatamente a mesma perda de treinamento podem apresentar capacidades de generalização completamente diferentes?}
\end{center}

A resposta está relacionada não apenas ao valor mínimo encontrado, mas também à geometria local da função de perda ao redor desse mínimo.

\subsection{Paisagem de Perda}

Pode-se imaginar a função de perda como uma superfície de alta dimensão denominada \textbf{paisagem de perda} (\textit{loss landscape}).

Cada ponto dessa superfície corresponde a um conjunto específico de parâmetros da rede neural.

A altura da superfície representa o valor da função de perda.

Durante o treinamento, algoritmos de otimização, como o Gradiente Descendente Estocástico (SGD), percorrem essa superfície procurando regiões onde a perda seja mínima.

Embora essa ideia seja simples em duas dimensões, nas redes neurais modernas a paisagem de perda é definida em espaços com milhões de dimensões, tornando impossível sua visualização direta.

Ainda assim, diversas propriedades geométricas dessa superfície podem ser estudadas matematicamente.

Entre elas, destaca-se a curvatura local da função de perda.

\subsection{Mínimos Sharp e Flat}

Nem todos os mínimos encontrados durante a otimização possuem as mesmas características geométricas.

Alguns mínimos apresentam curvatura extremamente elevada.

Nesses casos, pequenas alterações nos parâmetros produzem grandes aumentos no valor da perda.

Esses mínimos recebem o nome de \textbf{sharp minima}.

Por outro lado, existem mínimos cuja vizinhança permanece praticamente constante.

Mesmo após pequenas perturbações nos parâmetros, o valor da perda sofre pouca alteração.

Esses são denominados \textbf{flat minima}.

Diversos estudos experimentais indicam que modelos que convergem para mínimos planos tendem a apresentar melhor capacidade de generalização, pois pequenas perturbações inevitáveis durante a inferência não comprometem significativamente o desempenho do modelo.

Em contraste, modelos localizados em mínimos muito estreitos costumam ser altamente sensíveis a pequenas variações dos parâmetros.

\subsection{A Hessiana como Medida de Curvatura}

A principal ferramenta matemática utilizada para quantificar essa curvatura é a matriz Hessiana.

Seja

\[
\mathcal{L}(\theta)
\]

uma função de perda duas vezes continuamente diferenciável.

Sua Hessiana é definida como

\[
H(\theta)
=
\nabla_{\theta}^{2}\mathcal{L}(\theta),
\]

isto é, a matriz composta por todas as derivadas parciais de segunda ordem da função.

Enquanto o gradiente informa a direção de maior crescimento da perda, a Hessiana descreve como essa direção varia localmente.

Em particular, seus autovalores medem a curvatura da superfície ao longo de diferentes direções do espaço de parâmetros.

O maior autovalor da Hessiana,

\[
\lambda_{\max},
\]

corresponde justamente à direção de maior curvatura da paisagem de perda.

Quanto maior esse valor, mais "estreito" tende a ser o mínimo encontrado.

Assim, estimar o maior autovalor da Hessiana fornece uma medida quantitativa da nitidez (\textit{sharpness}) do mínimo.

\subsection{Desafio Computacional}

Apesar de sua importância teórica, calcular explicitamente a Hessiana é praticamente impossível em redes neurais modernas.

Se uma arquitetura possui

\[
n
\]

parâmetros, então a Hessiana possui dimensão

\[
n \times n,
\]

o que implica custo de armazenamento da ordem de

\[
\mathcal{O}(n^{2})
\]

e custo computacional para decomposição espectral da ordem de

\[
\mathcal{O}(n^{3}).
\]

Para redes com centenas de milhões de parâmetros, a construção explícita dessa matriz torna-se inviável tanto em memória quanto em tempo de processamento.

Surge então uma nova questão:

\begin{center}
\textit{Como estimar o maior autovalor da Hessiana sem nunca construir a Hessiana?}
\end{center}

A resposta envolve duas ferramentas fundamentais da Álgebra Linear e da Diferenciação Automática:

\begin{itemize}
    \item o Método da Potência, responsável por aproximar iterativamente o maior autovalor de uma matriz utilizando apenas produtos matriz-vetor;
    \item o Truque de Pearlmutter, que permite calcular exatamente o produto Hessiana-vetor sem construir explicitamente a Hessiana.
\end{itemize}

A combinação dessas técnicas possibilita estimar propriedades espectrais da Hessiana mesmo em redes neurais com milhões de parâmetros.

Por fim, essas informações são diretamente relacionadas ao algoritmo \textit{Sharpness-Aware Minimization} (SAM), cujo objetivo consiste em conduzir a otimização para mínimos geometricamente mais planos, favorecendo uma melhor capacidade de generalização.

%%%%%%%%%%%%%%%%%%%%%%%%%%%%%%%%%%%%%%%%%%%%%%%%%%%%%%%%%%%%%%%%%%%%%%%%%%
\section{Fundamentos de Álgebra Linear}

Antes de estudar o Método da Potência, o Truque de Pearlmutter e o algoritmo
Sharpness-Aware Minimization (SAM), é importante revisar alguns conceitos de
Álgebra Linear que aparecem continuamente ao longo deste trabalho.

Embora esses conceitos sejam amplamente conhecidos, compreender sua
interpretação geométrica facilita significativamente o entendimento dos
algoritmos apresentados nas seções seguintes.

%%%%%%%%%%%%%%%%%%%%%%%%%%%%%%%%%%%%%%%%%%%%%%%%%%%%%%%%%%%%%%%%%%%%%%%%%%
\subsection{Vetores como Pontos e Direções}

Em Álgebra Linear, um vetor pode ser interpretado de duas maneiras
equivalentes.

A primeira consiste em enxergá-lo como um ponto no espaço.

Por exemplo,

\[
x=
\begin{bmatrix}
2\\
1
\end{bmatrix}
\]

representa o ponto de coordenadas $(2,1)$ no plano.

Entretanto, na maior parte deste trabalho será utilizada uma segunda
interpretação, muito mais importante para Aprendizado de Máquina.

Um vetor será visto como uma direção associada a uma determinada magnitude.

Por exemplo,

\[
v=
\begin{bmatrix}
3\\
4
\end{bmatrix}
\]

indica uma direção específica do plano cujo comprimento é

\[
\|v\|_2
=
\sqrt{3^2+4^2}
=
5.
\]

Essa interpretação será fundamental quando estudarmos gradientes,
autovetores e produtos Hessiana-vetor.

\paragraph{Exemplo.}

Durante o treinamento de uma rede neural, o gradiente

\[
\nabla L(\theta)
\]

é um vetor.

Ele não representa um ponto da função de perda.

Em vez disso, indica a direção na qual a perda cresce mais rapidamente.

Logo, atualizar os parâmetros utilizando

\[
\theta_{k+1}
=
\theta_k-\eta\nabla L(\theta_k)
\]

significa mover-se na direção oposta ao gradiente.

%%%%%%%%%%%%%%%%%%%%%%%%%%%%%%%%%%%%%%%%%%%%%%%%%%%%%%%%%%%%%%%%%%%%%%%%%%
\subsection{Normas}

Em diversas partes deste trabalho será necessário medir o tamanho de um vetor.

Essa medida recebe o nome de norma.

A norma Euclidiana é definida por

\[
\|x\|_2
=
\sqrt{\sum_{i=1}^{n}x_i^2}.
\]

Ela corresponde exatamente à distância do vetor até a origem.

\paragraph{Exemplo.}

Considere

\[
x=
\begin{bmatrix}
1\\
2\\
2
\end{bmatrix}.
\]

Temos

\[
\|x\|_2
=
\sqrt{1+4+4}
=
3.
\]

No algoritmo SAM, essa norma aparece ao normalizar o gradiente,

\[
\hat{\epsilon}
=
\rho
\frac{\nabla L(w)}
{\|\nabla L(w)\|}.
\]

Observe que apenas a direção do gradiente é utilizada.

Seu comprimento é removido pela normalização.

%%%%%%%%%%%%%%%%%%%%%%%%%%%%%%%%%%%%%%%%%%%%%%%%%%%%%%%%%%%%%%%%%%%%%%%%%%
\subsection{Transformações Lineares}

Uma matriz pode ser interpretada como uma transformação aplicada aos vetores.

Considere

\[
A=
\begin{bmatrix}
2&0\\
0&1
\end{bmatrix}.
\]

Aplicando essa matriz ao vetor

\[
x=
\begin{bmatrix}
1\\
2
\end{bmatrix},
\]

obtém-se

\[
Ax=
\begin{bmatrix}
2\\
2
\end{bmatrix}.
\]

Observe que apenas a coordenada horizontal foi ampliada.

Portanto, multiplicar um vetor por uma matriz significa aplicar uma
transformação geométrica.

Essa interpretação será essencial para compreender o Método da Potência.

%%%%%%%%%%%%%%%%%%%%%%%%%%%%%%%%%%%%%%%%%%%%%%%%%%%%%%%%%%%%%%%%%%%%%%%%%%
\subsection{Autovalores e Autovetores}

Em geral, uma transformação linear modifica completamente um vetor.

Entretanto, existem direções especiais que permanecem inalteradas,
sofrendo apenas um fator de escala.

Essas direções recebem o nome de autovetores.

Formalmente,

\[
Av=\lambda v.
\]

O vetor

\[
v
\]

é chamado autovetor, enquanto

\[
\lambda
\]

é denominado autovalor.

O autovalor informa quanto aquela direção foi expandida ou contraída.

\paragraph{Exemplo.}

Considere

\[
A=
\begin{bmatrix}
3&0\\
0&1
\end{bmatrix}.
\]

Aplicando a transformação sobre

\[
v_1=
\begin{bmatrix}
1\\
0
\end{bmatrix},
\]

temos

\[
Av_1=
\begin{bmatrix}
3\\
0
\end{bmatrix}
=
3v_1.
\]

Logo,

\[
\lambda_1=3.
\]

Da mesma forma,

\[
v_2=
\begin{bmatrix}
0\\
1
\end{bmatrix}
\]

satisfaz

\[
Av_2=v_2,
\]

portanto

\[
\lambda_2=1.
\]

Observe que a transformação estica muito mais a direção horizontal do que
a vertical.

%%%%%%%%%%%%%%%%%%%%%%%%%%%%%%%%%%%%%%%%%%%%%%%%%%%%%%%%%%%%%%%%%%%%%%%%%%
\subsection{Interpretação Geométrica dos Autovalores}

Os autovalores medem a intensidade com que uma transformação linear altera
cada direção do espaço.

Quanto maior o módulo do autovalor,

\[
|\lambda|,
\]

maior será o alongamento daquela direção.

No Método da Potência veremos que, após sucessivas multiplicações pela matriz,

\[
A^k x,
\]

a componente associada ao maior autovalor cresce exponencialmente mais
rápido que todas as demais.

Esse comportamento explica por que o algoritmo converge para o autovetor
dominante.

%%%%%%%%%%%%%%%%%%%%%%%%%%%%%%%%%%%%%%%%%%%%%%%%%%%%%%%%%%%%%%%%%%%%%%%%%%
\subsection{Quociente de Rayleigh}

Depois que um vetor se aproxima do autovetor dominante, torna-se possível
estimar seu autovalor através do Quociente de Rayleigh,

\[
R(x)
=
\frac{x^TAx}{x^Tx}.
\]

Quando

\[
x
\]

é exatamente um autovetor,

\[
R(x)=\lambda.
\]

Na prática, mesmo antes da convergência completa, esse quociente fornece
excelentes aproximações do maior autovalor.

Essa propriedade será utilizada diretamente no Método da Potência para
estimar o maior autovalor da Hessiana.

%%%%%%%%%%%%%%%%%%%%%%%%%%%%%%%%%%%%%%%%%%%%%%%%%%%%%%%%%%%%%%%%%%%%%%%%%%
\subsection*{Resumo}

Nesta seção foram revisados os principais conceitos de Álgebra Linear
necessários para compreender o restante do trabalho.

Em particular, deve-se lembrar que

\begin{itemize}
\item vetores podem representar direções;
\item normas medem o comprimento de vetores;
\item matrizes representam transformações lineares;
\item autovetores são direções preservadas por uma transformação;
\item autovalores indicam o fator de expansão dessas direções;
\item o Quociente de Rayleigh fornece uma estimativa do autovalor associado a
um vetor.
\end{itemize}

Esses conceitos serão utilizados continuamente nas próximas seções para
explicar o Método da Potência, o Truque de Pearlmutter e o algoritmo SAM.




\section{Conceitos Utilizados no Projeto}

Esta seção conecta os conceitos matemáticos discutidos anteriormente ao
escopo específico deste projeto. A ideia central é interpretar o treinamento
de uma rede neural pequena, aplicada a bases como Iris e Titanic, como um
problema geométrico em um espaço de parâmetros.

Mesmo quando o modelo utilizado é pequeno em comparação com redes profundas
modernas, os mesmos fenômenos aparecem: a função de perda define uma paisagem,
o otimizador percorre essa paisagem, e o mínimo encontrado pode ser mais plano
ou mais agudo.

\subsection{Dados, Modelo e Parâmetros}

Nos experimentos do projeto, os exemplos de entrada vêm de bases tabulares.
No caso da base Iris, cada flor é descrita por medidas como comprimento e
largura da sépala e da pétala. No caso da base Titanic, cada passageiro é
descrito por atributos como classe da passagem, sexo, idade e outras
informações disponíveis no conjunto de dados.

Após o pré-processamento, cada exemplo pode ser representado por um vetor

\[
x_i \in \mathbb{R}^{d},
\]

em que \(d\) é o número de atributos usados pelo modelo. O rótulo correto é
representado por

\[
y_i.
\]

O modelo, por sua vez, possui um conjunto de pesos e vieses. Embora esses
parâmetros estejam distribuídos em várias camadas, é conveniente agrupá-los em
um único vetor:

\[
\theta =
\begin{bmatrix}
\theta_1 & \theta_2 & \cdots & \theta_n
\end{bmatrix}^{T}.
\]

Essa representação é exatamente a usada no código quando os parâmetros do
modelo são ``achatados'' em um vetor unidimensional. Assim, uma rede neural
pode ser estudada como um ponto em \(\mathbb{R}^{n}\).

\paragraph{Exemplo no contexto do projeto.}

Suponha uma MLP simples para o Titanic com \(d=6\) atributos de entrada, uma
camada escondida com \(10\) neurônios e uma saída binária. Apenas a primeira
camada já possui \(6 \times 10 = 60\) pesos, além de \(10\) vieses. A camada de
saída possui mais \(10\) pesos e \(1\) viés. Portanto, mesmo uma rede muito
pequena possui dezenas de parâmetros.

Cada configuração possível desses parâmetros corresponde a um ponto diferente
na paisagem de perda.

\subsection{Perda Empírica}

A perda empírica mede o erro médio do modelo sobre o conjunto de treinamento.
Se o conjunto de treino possui \(m\) exemplos, ela pode ser escrita como

\[
L_{\text{train}}(\theta)
=
\frac{1}{m}
\sum_{i=1}^{m}
\ell(f_{\theta}(x_i), y_i),
\]

em que \(f_{\theta}(x_i)\) é a predição do modelo e \(\ell\) é a função de
perda aplicada a um único exemplo.

Para classificação binária, como no Titanic, uma escolha comum é a entropia
cruzada binária. Para classificação multiclasse, como no Iris, utiliza-se
frequentemente a entropia cruzada categórica.

\paragraph{Interpretação.}

Minimizar \(L_{\text{train}}\) significa procurar parâmetros que expliquem bem
os exemplos de treino. Porém, o objetivo real do aprendizado de máquina é obter
baixa perda também em exemplos não vistos. Por isso, acompanhamos também uma
perda de validação ou teste:

\[
L_{\text{test}}(\theta).
\]

A diferença entre essas duas quantidades é uma forma simples de medir o
quanto o modelo está generalizando:

\[
\text{gap de generalização}
=
L_{\text{test}}(\theta)
-
L_{\text{train}}(\theta).
\]

Um gap pequeno sugere que o modelo aprendeu padrões que aparecem fora do
conjunto de treino. Um gap grande sugere que o modelo pode ter se ajustado
demais aos exemplos observados.

\subsection{Gradiente como Direção de Atualização}

Durante o treinamento, o gradiente da perda em relação aos parâmetros indica a
direção de maior aumento local da perda:

\[
\nabla_{\theta} L(\theta).
\]

Por isso, métodos como SGD e Adam atualizam os parâmetros no sentido oposto:

\[
\theta_{t+1}
=
\theta_t
-
\eta \nabla_{\theta} L(\theta_t),
\]

em que \(\eta\) é a taxa de aprendizado.

Geometricamente, o otimizador está caminhando sobre a paisagem de perda. Em
regiões inclinadas, o gradiente possui norma maior. Perto de um mínimo, o
gradiente tende a ficar pequeno. Entretanto, um gradiente pequeno não informa
sozinho se o mínimo é plano ou agudo. Para isso, precisamos olhar para a
curvatura.

\subsection{Hessiana e Aproximação Local da Loss}

A Hessiana descreve como a perda se curva ao redor de um ponto. Uma maneira
importante de entender seu papel é usar a expansão de Taylor de segunda ordem:

\[
L(\theta + \epsilon)
\approx
L(\theta)
+
\nabla L(\theta)^T \epsilon
+
\frac{1}{2}\epsilon^T H(\theta)\epsilon.
\]

Perto de um mínimo local, o gradiente \(\nabla L(\theta)\) é pequeno. Assim, o
termo que mais explica o aumento da perda após uma pequena perturbação
\(\epsilon\) é

\[
\frac{1}{2}\epsilon^T H(\theta)\epsilon.
\]

Esse termo mostra por que a Hessiana é tão relevante para o projeto: ela
quantifica quanto a perda aumenta quando os parâmetros são levemente
perturbados.

\paragraph{Exemplo numérico simples.}

Considere duas funções quadráticas:

\[
L_1(x,y)=10x^2+y^2
\quad \text{e} \quad
L_2(x,y)=x^2+y^2.
\]

Ambas possuem mínimo em \((0,0)\). Porém, \(L_1\) cresce muito mais rapidamente
na direção \(x\). Suas Hessianas são

\[
H_1=
\begin{bmatrix}
20&0\\
0&2
\end{bmatrix},
\quad
H_2=
\begin{bmatrix}
2&0\\
0&2
\end{bmatrix}.
\]

O maior autovalor de \(H_1\) é \(20\), enquanto o maior autovalor de \(H_2\) é
\(2\). Portanto, \(L_1\) possui um mínimo mais agudo, mesmo que as duas funções
tenham o mesmo valor mínimo.

Essa é a mesma ideia usada ao analisar uma rede neural treinada: não basta
olhar para o valor da loss no ponto final; é necessário observar como a loss
se comporta ao redor desse ponto.

\subsection{Sharpness e Perturbações nos Pesos}

No contexto deste projeto, sharpness pode ser entendida como a sensibilidade da
perda a pequenas mudanças nos parâmetros. Uma forma intuitiva de medir essa
sensibilidade é comparar

\[
L(\theta)
\quad \text{com} \quad
L(\theta + \epsilon),
\]

para perturbações pequenas \(\epsilon\).

Se pequenas perturbações causam grande aumento da loss, o mínimo é considerado
agudo. Se a loss quase não muda, o mínimo é considerado plano.

Essa ideia aparece diretamente nas visualizações implementadas no projeto. O
corte unidimensional avalia

\[
L(\theta^{*} + \alpha \delta),
\]

em que \(\theta^{*}\) representa os parâmetros treinados, \(\delta\) é uma
direção de perturbação e \(\alpha\) controla a intensidade dessa perturbação.

Já a visualização bidimensional avalia

\[
L(\theta^{*} + \alpha \delta_1 + \beta \delta_2),
\]

permitindo observar um pequeno pedaço da paisagem de perda ao redor da solução
encontrada.

\paragraph{Como ler os gráficos.}

Em um mínimo plano, o gráfico de \(L(\theta^{*}+\alpha\delta)\) tende a formar
um vale largo: a loss cresce lentamente quando \(\alpha\) se afasta de zero.
Em um mínimo agudo, a curva tende a subir rapidamente perto de \(\alpha=0\).

Nos mapas 2D, mínimos planos aparecem como regiões baixas e largas. Mínimos
agudos aparecem como poços estreitos, nos quais pequenas mudanças em
\(\alpha\) ou \(\beta\) já elevam bastante a loss.

\section{Método da Potência}

Após compreender os conceitos fundamentais de Álgebra Linear, podemos estudar um dos algoritmos numéricos mais importantes para problemas de grande escala: o Método da Potência (\textit{Power Iteration}).

Embora seja um algoritmo extremamente simples, ele desempenha um papel fundamental em diversas áreas da Computação, incluindo recuperação de informação (PageRank), processamento de sinais, visão computacional, física computacional e aprendizado profundo.

No contexto deste trabalho, seu objetivo consiste em estimar o maior autovalor da matriz Hessiana sem realizar sua decomposição espectral completa.

Essa estimativa será posteriormente utilizada para analisar a geometria da função de perda e compreender o funcionamento do algoritmo Sharpness-Aware Minimization (SAM).

Antes de apresentar sua formulação matemática, é importante compreender intuitivamente por que o algoritmo funciona.

\subsection{Intuição}

Imagine que uma matriz seja uma máquina capaz de deformar vetores.

Algumas direções são pouco alteradas.

Outras sofrem rotações.

Algumas são contraídas.

Entretanto, existe uma direção especial que sofre a maior ampliação possível.

Essa direção corresponde exatamente ao autovetor associado ao maior autovalor da matriz.

O Método da Potência explora justamente esse comportamento.

Em vez de calcular diretamente todos os autovalores — uma tarefa extremamente cara computacionalmente — o algoritmo aplica repetidamente a transformação linear sobre um vetor qualquer.

Como a direção associada ao maior autovalor cresce exponencialmente mais rápido que todas as demais, ela passa gradualmente a dominar o vetor resultante.

Após um número suficiente de iterações, praticamente toda a informação referente aos demais autovetores desaparece.

Consequentemente, o vetor obtido torna-se uma excelente aproximação do autovetor dominante.

\subsection{Exemplo Intuitivo}

Considere a matriz

\[
A=
\begin{bmatrix}
3&0\\
0&1
\end{bmatrix}.
\]

Essa matriz estica três vezes a direção horizontal, enquanto mantém inalterada a direção vertical.

Escolha um vetor inicial qualquer,

\[
x_0=
\begin{bmatrix}
1\\
1
\end{bmatrix}.
\]

Aplicando sucessivamente a transformação linear, obtém-se

\[
Ax_0=
\begin{bmatrix}
3\\
1
\end{bmatrix},
\]

\[
A^2x_0=
\begin{bmatrix}
9\\
1
\end{bmatrix},
\]

\[
A^3x_0=
\begin{bmatrix}
27\\
1
\end{bmatrix}.
\]

Observe que apenas a componente associada ao maior autovalor cresce rapidamente.

Após normalizações sucessivas,

\[
\frac{A^kx_0}{\|A^kx_0\|}
\longrightarrow
\begin{bmatrix}
1\\
0
\end{bmatrix},
\]

que corresponde exatamente ao autovetor dominante.

Esse pequeno exemplo ilustra praticamente toda a ideia por trás do Método da Potência.

Como os autovetores de uma matriz diagonalizável formam uma base para o espaço vetorial, qualquer vetor inicial pode ser escrito como combinação linear desses autovetores.

Essa decomposição é extremamente importante, pois permite analisar separadamente como cada componente evolui durante as sucessivas multiplicações pela matriz.

Assim,

\[
x_0
=
c_1v_1+c_2v_2+\cdots+c_nv_n.
\]

Cada coeficiente \(c_i\) mede quanto da direção \(v_i\) está presente no vetor inicial.

Caso \(c_1=0\), isto é, caso o vetor inicial seja exatamente ortogonal ao autovetor dominante, o algoritmo não consegue recuperá-lo.

Felizmente, essa situação possui probabilidade praticamente nula quando o vetor inicial é escolhido aleatoriamente.

\subsection{Formulação Algorítmica}

O Método da Potência pode ser descrito por três operações repetidas: multiplicar
pela matriz, normalizar o vetor resultante e estimar o autovalor associado.

Dada uma matriz \(A\) e um vetor inicial normalizado \(v_0\), cada iteração
calcula

\[
z_k = Av_k,
\]

seguida da normalização

\[
v_{k+1} = \frac{z_k}{\|z_k\|_2}.
\]

Ao final, o autovalor dominante é estimado pelo Quociente de Rayleigh:

\[
\lambda_{\max}
\approx
\frac{v_k^T A v_k}{v_k^T v_k}.
\]

No caso deste projeto, a matriz \(A\) é a Hessiana \(H\) da função de perda no
ponto encontrado após o treinamento. Portanto, a iteração conceitual passa a
ser

\[
z_k = H v_k.
\]

O ponto crucial é que o código nunca precisa construir \(H\). Ele precisa
apenas ser capaz de calcular o produto \(Hv_k\), tarefa realizada pelo Truque
de Pearlmutter.

\begin{algorithm}[H]
\DontPrintSemicolon
\Input{Parâmetros treinados \(\theta^{*}\), perda \(L\), número de iterações \(T\)}
\Output{Estimativa de \(\lambda_{\max}(H)\)}
\BlankLine
Escolha um vetor aleatório \(v_0\) com \(\|v_0\|_2=1\)\;
\Para{\(k = 0, 1, \dots, T-1\)}{
    Calcule \(z_k \leftarrow H(\theta^{*})v_k\) usando produto Hessiana-vetor\;
    Estime \(\lambda_k \leftarrow v_k^T z_k\)\;
    Normalize \(v_{k+1} \leftarrow z_k / \|z_k\|_2\)\;
}
\Retorna{\(\lambda_T\)}\;
\caption{Método da Potência aplicado à Hessiana implícita}
\label{alg:power_hessian}
\end{algorithm}

\paragraph{Leitura no contexto do experimento.}

Quando o valor estimado de \(\lambda_{\max}(H)\) é alto, existe pelo menos uma
direção no espaço de parâmetros em que a loss cresce rapidamente. Essa direção
é justamente aproximada pelo autovetor dominante produzido pelo Método da
Potência. No experimento de perturbação direcional, perturbar os pesos nessa
direção deve causar um aumento maior de loss do que perturbar em direções
aleatórias comuns.

Quando \(\lambda_{\max}(H)\) é baixo, a maior curvatura local encontrada também
é pequena. Isso sugere que a região ao redor dos pesos treinados é mais estável
e, portanto, mais compatível com a noção geométrica de mínimo plano.

\subsection{Truque de Pearlmutter}
Em problemas de otimização de redes neurais profundas, o vetor de parâmetros possui dimensionalidade colossal ($n \gg 10^6$). O \textbf{Truque de Pearlmutter} (ou produto Hessiana-vetor exato) resolve a limitação de memória ao possibilitar o cálculo direto do produto $H v$, onde $v \in \mathbb{R}^n$ é um vetor arbitrário, sob complexidade de tempo e espaço de ordem \textbf{$\mathcal{O}(n)$}, tornando-o computacionalmente equivalente ao custo de um gradiente tradicional.

\subsubsection{Formulação Matemática e o Operador $\mathcal{R}\{\cdot\}$}
Seja $\theta \in \mathbb{R}^n$ o vetor de parâmetros da rede neural e $\mathcal{L}(\theta)$ a função de perda. O gradiente é definido por $g(\theta) = \nabla_{\theta} \mathcal{L}(\theta) \in \mathbb{R}^n$ e a matriz Hessiana por $H(\theta) = \nabla_{\theta}^2 \mathcal{L}(\theta) \in \mathbb{R}^{n \times n}$. Objetiva-se computar:
\begin{equation}
w = H(\theta) v
\end{equation}

Pearlmutter introduziu o operador $\mathcal{R}_v\{\cdot\}$, que define a derivada direcional de uma função de parâmetros na direção do vetor $v$, avaliada sob uma perturbação escalar infinitesimal $e$:
\begin{equation}
\mathcal{R}_v \{ f(\theta) \} \equiv \left. \frac{d}{de} f(\theta + ev) \right|_{e=0}
\end{equation}

Aplicando o operador $\mathcal{R}_v\{\cdot\}$ diretamente sobre o vetor do gradiente $g(\theta)$, obtém-se via Regra da Cadeia multivariável:
\begin{equation}
\mathcal{R}_v \{ g(\theta) \} = \left. \frac{d}{de} \nabla_{\theta} \mathcal{L}(\theta + ev) \right|_{e=0} = \nabla_{\theta}^2 \mathcal{L}(\theta) \cdot v = H(\theta)v
\end{equation}

Esta igualdade demonstra analiticamente que o produto Hessiana-vetor pode ser obtido aplicando uma diferenciação unidimensional comum sobre o operador de primeira ordem da rede.

\subsubsection{Implementação Computacional via Diferenciação Automática}
Nos arcabouços modernos de computação gráfica e diferenciação automática (como PyTorch e TensorFlow), o Truque de Pearlmutter é executado de forma exata através de duas etapas principais:
\begin{enumerate}
    \item \textbf{Passada Forward e Backward Inicial:} Avalia-se a função de perda $\mathcal{L}(\theta)$ e computa-se o gradiente convencional via retropropagação padrão: $g(\theta) = \nabla_{\theta} \mathcal{L}(\theta)$.
    \item \textbf{Projeção Escalar e Segunda Passada:} Realiza-se o produto escalar entre o gradiente obtido e o vetor arbitrário $v$, gerando um escalar funcional $\phi$:
    \begin{equation}
    \phi = g(\theta)^T v = \sum_{i=1}^n \frac{\partial \mathcal{L}}{\partial \theta_i} v_i
    \end{equation}
    Calcula-se o gradiente deste novo escalar $\phi$ em relação aos parâmetros originais $\theta$. Como $v$ é tratado como uma constante nesta etapa, a linearidade do operador de derivada resulta em:
    \begin{equation}
    \nabla_{\theta} \phi = \nabla_{\theta} (g(\theta)^T v) = \nabla_{\theta}^2 \mathcal{L}(\theta) v = H v
    \end{equation}
\end{enumerate}

---

\subsection{SAM (Sharpness-Aware Minimization)}
O algoritmo \textit{Sharpness-Aware Minimization} (SAM) é um procedimento de otimização desenvolvido para aprimorar a capacidade de generalização ao modificar o objetivo de treinamento. Em vez de focar estritamente em um ponto isolado de perda mínima nos dados de treino, o SAM busca uma região no espaço de parâmetros onde toda a vizinhança apresente, uniformemente, valores baixos de perda. Geometricamente, esta formulação força o modelo a convergir para mínimos planos (\textit{flat minima}) em detrimento de mínimos estreitos (\textit{sharp minima}), mitigando cenários de \textit{overfitting}.

\subsubsection{Formulação Matemática e o Problema Min-Max}
Seja $\mathcal{S} \triangleq \bigcup_{i=1}^{n} \{(x_i, y_i)\}$ o conjunto de dados de treinamento amostrado de uma distribuição $\mathcal{D}$. A perda empírica é dada por $L_{\mathcal{S}}(w) \triangleq \frac{1}{n}\sum_{i=1}^{n} l(w, x_i, y_i)$, onde $w \in \mathbb{R}^d$ representa o vetor de pesos. Baseado em limites teóricos de generalização via formalismo PAC-Bayesiano, o SAM propõe o seguinte problema de otimização min-max:
\begin{equation}
\min_{w} L_{\mathcal{S}}^{\text{SAM}}(w) + \lambda \|w\|_2^2 \quad \text{onde} \quad L_{\mathcal{S}}^{\text{SAM}}(w) \triangleq \max_{\|\epsilon\|_2 \le \rho} L_{\mathcal{S}}(w + \epsilon)
\end{equation}
onde $\rho > 0$ define o raio da vizinhança esférica sob a norma $\ell_2$, e $\lambda$ controla a regularização por decaimento de pesos (\textit{weight decay}). A métrica de nitidez (\textit{sharpness}) implícita avalia o crescimento local da perda:
\begin{equation}
\text{Sharpness}(w) = \max_{\|\epsilon\|_2 \le \rho} L_{\mathcal{S}}(w + \epsilon) - L_{\mathcal{S}}(w)
\end{equation}

\subsubsection{Linearização do Problema Interno e Vetor Adversário}
Para viabilizar computacionalmente a maximização interna, aplica-se uma expansão em série de Taylor de primeira ordem para $L_{\mathcal{S}}(w+\epsilon)$ em torno de $\epsilon = 0$:
\begin{equation}
\epsilon^*(w) \approx \arg\max_{\|\epsilon\|_2 \le \rho} \left[ L_{\mathcal{S}}(w) + \epsilon^T \nabla_w L_{\mathcal{S}}(w) \right] = \arg\max_{\|\epsilon\|_2 \le \rho} \epsilon^T \nabla_w L_{\mathcal{S}}(w)
\end{equation}

A solução analítica para este problema sob restrição euclidiana decorre diretamente da aplicação da desigualdade de Cauchy-Schwarz, onde a projeção máxima ocorre quando $\epsilon$ orienta-se colinearmene ao vetor gradiente $\nabla_w L_{\mathcal{S}}(w)$. Assim, a perturbação ótima escalada assume a forma:
\begin{equation}
\hat{\epsilon}(w) = \rho \frac{\nabla_w L_{\mathcal{S}}(w)}{\|\nabla_w L_{\mathcal{S}}(w)\|_2}
\end{equation}

\subsubsection{Aproximação do Gradiente de SAM e Termos de Segunda Ordem}
Substituindo a perturbação ótima $\hat{\epsilon}(w)$ de volta no objetivo principal e diferenciando em relação a $w$, obtém-se o gradiente efetivo via regra da cadeia:
\begin{equation}
\nabla_w L_{\mathcal{S}}^{\text{SAM}}(w) \approx \nabla_w L_{\mathcal{S}}(w)\big|_{w + \hat{\epsilon}(w)} + \left( \frac{d\hat{\epsilon}(w)}{dw} \right) \nabla_w L_{\mathcal{S}}(w)\big|_{w + \hat{\epsilon}(w)}
\end{equation}

O termo diferencial $\frac{d\hat{\epsilon}(w)}{dw}$ introduz a matriz Hessiana de segunda ordem na equação. Embora passível de cálculo exato via Truque de Pearlmutter, constatou-se empiricamente em \cite{SAM} que o descarte dos termos de segunda ordem acelera consideravelmente o tempo de processamento por época e preserva as propriedades matemáticas de generalização. A aproximação final simplificada do gradiente do SAM estabelece-se como:
\begin{equation}
\nabla_w L_{\mathcal{S}}^{\text{SAM}}(w) \approx \nabla_w L_{\mathcal{S}}(w)\big|_{w + \hat{\epsilon}(w)}
\end{equation}

\subsubsection{Dinâmica do Algoritmo}
A dinâmica operacional do SAM requer \textbf{duas passadas completas de retropropagação} por iteração sobre um determinado subconjunto de dados (\textit{mini-batch}) $\mathcal{B}$, conforme descrito formalmente no Algoritmo \ref{alg:sam}.

\begin{algorithm}[H]
\DontPrintSemicolon
\Input{Conjunto de treino $\mathcal{S}$, Taxa de aprendizado $\eta$, Raio da perturbação $\rho$, Pesos iniciais $w_0$}
\Output{Pesos otimizados $w$}
\BlankLine
\Para{cada iteração $t = 0, 1, 2, \dots$}{
    Amostre um mini-batch $\mathcal{B} \subset \mathcal{S}$\;
    Calcule o gradiente empírico inicial: $g_t \leftarrow \nabla_w L_{\mathcal{B}}(w_t)$\;
    Computar a perturbação adversária ótima: $\hat{\epsilon}_t \leftarrow \rho \frac{g_t}{\|g_t\|_2}$\;
    Calcular a coordenada de subida máxima: $w_{\text{adv}} \leftarrow w_t + \hat{\epsilon}_t$\;
    Avaliar o gradiente efetivo na posição perturbada: $g_{\text{SAM}} \leftarrow \nabla_w L_{\mathcal{B}}(w)\big|_{w_{\text{adv}}}$\;
    Atualizar o vetor de parâmetros original: $w_{t+1} \leftarrow w_t - \eta g_{\text{SAM}}$\;
}
\Retorna{$w$}\;
\caption{Algoritmo Sharpness-Aware Minimization (SAM)}
\label{alg:sam}
\end{algorithm}

\begin{thebibliography}{1}

\bibitem{SAM} FORET, Pierre; KLEINER, Ariel; MOBAHI, Hossein; NEYSHABUR, Behnam. Sharpness-aware minimization for efficiently improving generalization. In: \textbf{International Conference on Learning Representations (ICLR)}, 2021, Virtual. Anais... Virtual: ICLR, 2021.

\end{thebibliography}

\end{document}
